% 来源：图片838—840教材数学文化；按材料概述，不将现代log定义等同早期对数
\begin{frame}[t]{附录：对数为何被发明}
对数曾是简化繁重计算的重要工具。所附教材介绍：
\begin{itemize}
\item 纳皮尔于1614年发表对数著作。
\item 别尔基也独立研究了对数，相关著作于1620年出版。
\item 布里格斯推动常用对数表的编制。
\end{itemize}
\[2^6\cdot2^8=2^{6+8}=2^{14},\qquad64\times256=16384.\]
将幂对应到指数，就能把乘法转化为加法。\end{frame}

% 来源：图片838—839
\begin{frame}[t]{附录：从稀疏表到精细表}
\small
\centering\begin{tabular}{c|rrrrrrrr}
$n$&0&1&2&3&4&5&6&7\\\hline
$2^n$&1&2&4&8&16&32&64&128\\\hline
$n$&8&9&10&11&12&13&14&\\\hline
$2^n$&256&512&1024&2048&4096&8192&16384&
\end{tabular}\par\normalsize
表中缺少3、5、7等数，怎样把间隔缩小？

取接近1的 $q=1.0001$：每次把 $x$ 增加其万分之一，
\[x_0=1,\quad x_{n+1}=x_n+\frac{x_n}{10000}=1.0001x_n,\quad x_n=1.0001^n.\]\end{frame}

% 来源：图片839
\begin{frame}[t]{附录：历史表的近似数}
\[x=1.0001^y,\qquad y=\log_{1.0001}x.\]
\[1.0001^{100}\approx1.010049662,\quad
1.0001^{101}\approx1.010150667.\]
\[1.0001^{46054}\approx99.999955936.\]

教材用接近1的底数解释早期制表思想。数值表经舍入，只能得到近似计算结果；这些现代符号是对思想的解释。\end{frame}

% 来源：图片839—840
\begin{frame}[t]{附录：查表把除法变成减法}
\small\centering\renewcommand{\arraystretch}{1.35}
\begin{tabular}{c|rrrrr}
$y$&20333&20334&20335&20336&20337\\
$x\approx$&7.6385&7.6392&7.6400&7.6408&7.6415\\\hline
$y$&10237&10238&10239&10240&10241\\
$x\approx$&2.7833&2.7836&2.7839&2.7842&2.7844
\end{tabular}\par\normalsize
求 $7640\div2784=7.640\div2.784$。

按表取对应指数 $20335$ 和 $10239$，相减为 $10096$。\end{frame}

% 来源：图片840；修正原近似分母使用精确等号的问题
\begin{frame}[t]{附录：近似计算与换底}
\small\centering\begin{tabular}{c|rrrrr}
$y$&10094&10095&10096&10097&10098\\\hline
$x\approx$&2.7438&2.7441&2.7444&2.7446&2.7449
\end{tabular}\par\normalsize
查表得到商约为 $2.7444$；实际商为 $2.74425287\ldots$。
两者保留4位有效数字均为 $2.744$。

\[\lg x=\frac{\log_{1.0001}x}{\log_{1.0001}10}
\approx\frac{\log_{1.0001}x}{23027}.\]
本页换底公式作为后续学习的预告，分母 $23027$ 是近似值。\end{frame}

% 来源：图片841顶部
\begin{frame}[t]{补充：教材思考题}
由 $0.568^x=0.01$，可准确记为
\[x=\log_{0.568}0.01.\]
借助计算器得 $x\approx8.142$。

对数符号能够准确表示指数，近似小数方便实际计算，两者作用不同。\end{frame}
% 来源：图片841顶部；图片839的制表背景。
\begin{frame}[t]{补充：指数结构与制表背景}
\[2^{\log_23-2}=\frac{2^{\log_23}}{2^2}=\frac34.\]
\[2^{\log_2(3+2)}=5.\]
括号内的加法与指数上的减法，处理方式不同。

所附数学文化材料还介绍：纳皮尔采用接近1的小于1的数
$b=0.9999999$，别尔基采用 $a=1.0001$，以缩小表中相邻数的间隔。

现代指数记号和对数与指数的统一认识经历了长期发展，早期对数表不能直接等同于今天的常用对数表。
\end{frame}
